\chapter{Tài liệu}
\label{chuong:tai-lieu}

Thầy cô tải sách giáo khoa hay tài liệu riêng lên để Kriky ra đề bám đúng nội dung ấy.

\viec{Tải tài liệu lên}

\anh[0.5\linewidth]{T/13-tai-lieu}{
  \diem{1}{0.82}{0.13}
  \diem{2}{0.62}{0.52}
}

\begin{buoc}
  \item Ở nhóm \nhan{TÀI LIỆU} của thanh trái, bấm nút tải lên rồi chọn tệp.
  \item Đợi tài liệu đọc xong. Sách dày có thể mất vài phút.
\end{buoc}

\xong{dưới tên tệp không còn chữ \nhan{Đang xử lý…} hay \nhan{Đang đọc …}; tệp PDF có thêm số trang.}

\luuy{Kriky nhận tệp \go{.pdf}, \go{.txt} và \go{.md}, nặng tối đa 100 MB. Tệp Word chưa dùng được: hãy lưu
thành PDF trước.}

\viec{Đọc trạng thái của tài liệu}

\begin{bangtra}{L{4.6cm}Y}{\dau{Dưới tên tệp} & \dau{Nghĩa là, làm gì}}
  \nhan{Đang xử lý…} & Kriky bắt đầu đọc tệp. Đợi, không cần làm gì. \\
  \nhan{Đang đọc 12/104 trang…} & Đang đọc tới trang 12 trên 104. \\
  \nhan{Đang dựng mục lục…} & Đọc xong các trang, đang dựng mục lục. Sắp xong. \\
  Số trang, ví dụ \nhan{104 trang} & Đã đọc xong, dùng được để ra đề. Tệp \go{.txt}, \go{.md} không có số trang. \\
  \nhan{Không đọc được chữ} & Tệp không có chữ đọc được, ví dụ ảnh chụp quá mờ. Tải bản rõ hơn. \\
  \nhan{Xử lí lỗi. Hãy tải lại} & Đọc hỏng giữa chừng. Tải lại tệp ấy. \\
\end{bangtra}

\machnho{Tải tài liệu lên chưa làm Kriky dùng nó ngay. Muốn ra đề theo một tài liệu, thầy cô gắn tài
liệu ấy vào lời nhắn (Chương~\ref{chuong:soan-de}).}
